\documentclass[11pt]{article}
\usepackage{amsmath,amssymb,bm}
\usepackage[margin=1in]{geometry}

\newcommand{\Ms}{\mathcal{M}}
\newcommand{\Ds}{\mathcal{D}}
\newcommand{\Cs}{\mathcal{C}}
\newcommand{\Ks}{\mathcal{K}}
\newcommand{\Da}{\mathrm{Da}}
\newcommand{\Bi}{\mathrm{Bi}}
\newcommand{\Pe}{\mathrm{Pe}}
\newcommand{\Rhat}{\widehat{R}}
\newcommand{\pp}{\phi_{p0}}

\title{Initial Model Derivation from Migrated Appendix}
\author{Migrated Appendix A--B source material, organized for repository review}
\date{\today}

\begin{document}
\maketitle

\noindent This file is the source of truth for long mathematical derivations. Markdown files should only summarize it.

\medskip
\noindent\textbf{Provenance note.}
This document mechanically copies Appendix A, ``Model Construction,'' and Appendix B, ``Nondimensionalization,'' from \texttt{paper/appendix/appendix.tex}. The copied source begins at the Model Construction section and stops before the Homogeneous Diagnostic Derivation section. Wrapper packages and macro definitions were added only so this derivation can compile as a standalone document. No scientific corrections, equation edits, or validation claims are made here.

\section{Model Construction}
\label{app:model}

We formulate a one-dimensional Lagrangian model for the coupled thermo-chemo-poroelastic dynamics of a thermoresponsive gel slab undergoing an exothermic reaction. The conservative variables are the local swelling ratio $J(X,t)$, the reactant content per reference volume $Jc$, and the temperature $T(X,t)$; the solvent chemical potential is used as an algebraic mixed variable.

\subsection{Kinematics}

Consider a gel slab of reference (as-prepared) half-thickness $H_0$. We adopt a material (Lagrangian) coordinate $X\in[0,H_0]$, with $X=0$ denoting the symmetry plane and $X=H_0$ the free surface. The current position of a material point is $x(X,t)$, and the local volume ratio is
\begin{equation}
  J(X,t) = \frac{\partial x}{\partial X} > 0.
  \label{eq:J_def}
\end{equation}
In one dimension, $J$ is the local swelling ratio. The polymer volume fraction is a dependent variable determined by mass conservation of the network:
\begin{equation}
  \phi(X,t) = \frac{\pp}{J(X,t)},
  \label{eq:phi_J}
\end{equation}
where $\pp$ is the polymer volume fraction in the reference (as-prepared) state; $\pp=0.15$ corresponds to a gel that is $85\%$ solvent by volume at preparation. The current thickness of the slab is
\begin{equation}
  h(t) = x(H_0,t) = \int_0^{H_0} J(X,t)\,dX.
  \label{eq:thickness}
\end{equation}

\subsection{Governing equations}

\subsubsection{Swelling--permeation equation}

The evolution of $J$ is governed by solvent conservation:
\begin{equation}
  \frac{\partial J}{\partial t} = -\frac{\partial Q_s}{\partial X},
  \label{eq:J_evol}
\end{equation}
where $Q_s$ is the solvent flux relative to the polymer network, driven by chemical-potential gradients:
\begin{equation}
  Q_s = -M_s(J,T)\,\frac{\partial \mu_s}{\partial X}.
  \label{eq:Qs}
\end{equation}
Here $M_s(J,T)$ is the permeation mobility and $\mu_s$ is the solvent chemical potential.

\subsubsection{Reactant transport equation}

The reactant concentration $c(X,t)$ (moles per unit current total volume) evolves according to
\begin{equation}
  \frac{\partial (Jc)}{\partial t}
  + \frac{\partial N_c}{\partial X}
  = -J\,R(c,T,J),
  \label{eq:c_evol}
\end{equation}
with the total reactant flux
\begin{equation}
  N_c = c\,Q_s - D_c(J,T)\,\frac{\partial c}{\partial X}.
  \label{eq:Nc}
\end{equation}
The first term represents advective transport by the solvent flow, and the second term Fickian diffusion within the pore fluid. The effective diffusivity $D_c(J,T)$ depends on the local microstructure via $J$.

\subsubsection{Heat equation}

The temperature field satisfies
\begin{multline}
  C_\mathrm{eff}(J)\,\frac{\partial T}{\partial t}
  = \frac{\partial}{\partial X}\!\left(K_T(J)\,\frac{\partial T}{\partial X}\right)
  - \frac{\partial}{\partial X}\!\left(c_p^\ell\,T\,Q_s\right)
  \\
  + (-\Delta H_r)\,J\,R(c,T,J),
  \label{eq:T_evol}
\end{multline}
where $C_\mathrm{eff}(J)$ is the effective volumetric heat capacity,
$K_T(J)$ the effective thermal conductivity, $c_p^\ell$ the solvent
heat capacity per unit volume, and $-\Delta H_r > 0$ the exothermic
reaction enthalpy.  The three right-hand-side terms represent thermal
conduction, enthalpy advection by the solvent flow, and reaction heat
generation, respectively.

\subsection{Free energy and chemical potential}

The Helmholtz free-energy density per unit reference volume is written as
\begin{equation}
  \Psi(J,T) = \Psi_\mathrm{el}(J) + \Psi_\mathrm{mix}(\phi,T)
  + \frac{\kappa_J}{2}\left(\frac{\partial J}{\partial X}\right)^2,
  \label{eq:Psi}
\end{equation}
where $\phi = \pp/J$, and the three contributions are the elastic
energy of the network, the Flory--Huggins mixing energy, and a
gradient-energy regularization.

For the polymer network we adopt a neo-Hookean elastic energy
appropriate for uniaxial swelling in the slab geometry:
\begin{equation}
  \Psi_\mathrm{el}(J) = \frac{G}{2}
  \bigl(J^{2} - 1 - 2\ln J\bigr),
  \label{eq:Psi_el}
\end{equation}
where $G$ is the shear modulus of the reference-state network.

The mixing free-energy density takes the standard Flory--Huggins
form~\cite{Flory1953}:
\begin{equation}
  \Psi_\mathrm{mix}(\phi,T)
  = \frac{R_g T}{v_s}
  \Big[(1-\phi)\ln(1-\phi) + \chi(T,\phi)\,\phi(1-\phi)\Big],
  \label{eq:Psi_mix}
\end{equation}
with $R_g$ the gas constant, $v_s$ the solvent molar volume, and
the Flory--Huggins interaction parameter
\begin{equation}
  \chi(T,\phi) = \chi_0(T) + \chi_1\,\phi.
  \label{eq:chi}
\end{equation}
The temperature-dependent part $\chi_0(T)$ encodes the LCST
behavior of the gel; $\chi_1$ captures the leading
concentration dependence.

The solvent chemical potential is obtained as
$\mu_s = v_s\,\partial\Psi/\partial(1-\phi)\cdot
\partial(1-\phi)/\partial J$~\cite{Hong2008}.  Evaluating this
derivative yields the mixing contribution
\begin{equation}
  \frac{\mu_\mathrm{mix}}{R_g T/v_s}
  = \ln(1-\phi) + \phi + \chi(T,\phi)\,\phi^2
  - \phi^2(1-\phi)\,\frac{\partial\chi}{\partial\phi},
  \label{eq:mu_mix}
\end{equation}
so that the linear form~\eqref{eq:chi} gives
$\partial\chi/\partial\phi=\chi_1$.  The shorter expression without
this final term is retained only as an effective-osmotic comparison
closure, not as a strict free-energy derivative.
and the elastic contribution
\begin{equation}
  \frac{\mu_\mathrm{el}}{R_g T/v_s}
  = \frac{G\,v_s}{R_g T}
  \bigl(J - J^{-1}\bigr).
  \label{eq:mu_el}
\end{equation}
The full chemical potential, including the gradient-energy
regularization, is
\begin{equation}
  \mu_s = \mu_\mathrm{mix}(J,T) + \mu_\mathrm{el}(J)
  - \kappa_J\,\frac{\partial^2 J}{\partial X^2}.
  \label{eq:mu_s}
\end{equation}

\subsection{Reaction kinetics}

The volumetric reaction rate takes the Arrhenius form modulated by a collapse-dependent accessibility factor:
\begin{equation}
  R(c,T,J) = k_0\,c^n\,\exp\!\left(-\frac{E_a}{R_g T}\right) a(J),
  \label{eq:R}
\end{equation}
where $k_0$ is the pre-exponential factor, $n$ the reaction order, and $E_a$ the activation energy.

The activity factor $a(J)$ models the suppression of reaction due to reduced free solvent volume upon collapse:
\begin{equation}
  a(J) = A_\mathrm{min}+(1-A_\mathrm{min})\,
  \bar p^{\,m_\mathrm{act}}, \qquad
  \bar p=\mathrm{clip}\!\left[
    \frac{1-\phi}{1-\phi_\mathrm{ref}},0,1\right].
  \label{eq:aJ}
\end{equation}
When the gel collapses ($J\to\pp$, $\phi\to 1$), the activity approaches
the residual floor $A_\mathrm{min}$.  This is a phenomenological
porosity/accessibility closure rather than a microscopic derivation.

The effective transport coefficients similarly depend on the microstructure:
\begin{align}
  D_c(J) &= D_0\left[D_\mathrm{min}+(1-D_\mathrm{min})
  \bar p^{\,m_\mathrm{diff}}\right], \label{eq:Dc} \\
  M_s(J) &= M_0\,\bar p^{\,m_\mathrm{mob}}, \label{eq:Ms}
\end{align}
with optional Mackie--Meares and free-volume forms used only for
comparison.  Thus solute diffusion and solvent permeation are
suppressed but not forced to vanish in the collapsed state.

\subsection{Boundary conditions}

At the symmetry plane $X=0$:
\begin{equation}
  Q_s = 0, \qquad N_c = 0, \qquad \frac{\partial T}{\partial X} = 0.
  \label{eq:bc_sym}
\end{equation}

At the free surface $X=H_0$:

\paragraph{Solvent exchange.}
A Robin-type condition couples the surface chemical potential to the bath:
\begin{equation}
  Q_s = L_\mu\!\left(\mu_s - \mu_\mathrm{bath}\right),
  \label{eq:bc_mu}
\end{equation}
where $L_\mu$ is a surface transfer coefficient. In the limit of rapid equilibration, $\mu_s = \mu_\mathrm{bath}$.

\paragraph{Reactant supply.}
\begin{equation}
  N_c = k_c\,(c - c_b),
  \label{eq:bc_c}
\end{equation}
where $c_b$ is the constant bath concentration.

\paragraph{Heat exchange.}
\begin{equation}
  -K_T\,\frac{\partial T}{\partial X} = h_T\,(T - T_\infty),
  \label{eq:bc_T}
\end{equation}
where $h_T$ is the surface heat-transfer coefficient and $T_\infty$ the ambient temperature.

\paragraph{Mechanical boundary.}
The free surface is traction-free. In the $J$-formulation this is implicitly satisfied by the chemical-potential balance.

%% ================================================================
%%  Appendix B: Nondimensionalization
%% ================================================================
\section{Nondimensionalization}
\label{app:nondim}

\subsection{Scaling}

We introduce the following dimensionless variables. Length is scaled by the reference thickness:
\begin{equation}
  X = H_0\,\tilde{x}, \qquad \tilde{x}\in[0,1].
\end{equation}
The swelling ratio $J$ is already dimensionless. Reactant concentration is normalized by the bath value:
\begin{equation}
  c = c_b\,u.
\end{equation}
Temperature is measured relative to ambient, scaled by the adiabatic reaction temperature rise:
\begin{equation}
  T = T_\infty + \Delta T_*\,\theta,
  \qquad
  \Delta T_* = \frac{(-\Delta H_r)\,c_b}{C_*},
\end{equation}
where $C_*$ is a representative volumetric heat capacity.

The chemical potential is scaled by the thermal mixing energy:
\begin{equation}
  \mu_s = \mu_*\,m, \qquad \mu_* = \frac{R_g T_\infty}{v_s}.
\end{equation}

Time is scaled by the swelling/permeation time:
\begin{equation}
  \tau_s = \frac{H_0^2}{M_0\,\mu_*},
\end{equation}
where $M_0$ is a reference mobility. The solvent and reactant flux scales follow:
\begin{equation}
  Q_s = \frac{H_0}{\tau_s}\,q,
  \qquad
  N_c = \frac{c_b\,H_0}{\tau_s}\,n.
\end{equation}

Material functions are factored as reference value times dimensionless function:
\begin{align}
  M_s &= M_0\,\Ms(J,\theta), &
  D_c &= D_0\,\Ds_\mathrm{cur}(J,\theta), \\
  C_\mathrm{eff} &= C_*\,\Cs(J), &
  K_T &= K_*\,\Ks_\mathrm{cur}(J).
\end{align}
In the reference-coordinate equations below, the baseline closure
uses effective reference-coordinate coefficients
$\Ds_\mathrm{ref}$, $\Ks_\mathrm{ref}$, and $\Cs_\mathrm{ref}$.
When the tabulated coefficients are interpreted as current-space
material properties, the metric-corrected comparison uses
$\Ds_\mathrm{ref}=\Ds_\mathrm{cur}/J$,
$\Ks_\mathrm{ref}=\Ks_\mathrm{cur}/J$, and
$\Cs_\mathrm{ref}=J\Cs_\mathrm{cur}$.

\subsection{Dimensionless reaction rate}

Defining the reference reaction rate
\begin{equation}
  R_0 = k_0\,c_b^{\,n}\,\exp\!\left(-\frac{E_a}{R_g T_\infty}\right),
\end{equation}
we write $R = R_0\,\Rhat(u,\theta,J)$ with
\begin{equation}
  \Rhat(u,\theta,J)
  = u^n\,a(J)\,
  \exp\!\left[\frac{\Gamma_A\,\theta}{1+\varepsilon_T\,\theta}\right],
  \label{eq:Rhat}
\end{equation}
where
\begin{equation}
  \varepsilon_T = \frac{\Delta T_*}{T_\infty},
  \qquad
  \Gamma_A = \frac{E_a\,\Delta T_*}{R_g\,T_\infty^2}.
  \label{eq:epsT_GammaA}
\end{equation}
When $\varepsilon_T \ll 1$, this simplifies to $\Rhat \approx u^n\,a(J)\,e^{\Gamma_A\theta}$.

\subsection{Dimensionless chemical potential}

The dimensionless chemical potential is
\begin{equation}
  m = m_\mathrm{mix}(J,\theta) + m_\mathrm{el}(J) - \ell^2\,\frac{\partial^2 J}{\partial \tilde{x}^2},
  \label{eq:m_nondim}
\end{equation}
with the interface parameter
\begin{equation}
  \ell^2 = \frac{\kappa_J}{\mu_*\,H_0^2}.
\end{equation}

The mixing part retains the Flory--Huggins form (neglecting the
$T/T_\infty\approx 1+\varepsilon_T\theta\approx 1$ prefactor, valid
when $\varepsilon_T\ll 1$):
\begin{equation}
  m_\mathrm{mix} = \ln(1-\phi) + \phi + \chi(\theta,\phi)\,\phi^2
  -\phi^2(1-\phi)\,\frac{\partial\chi}{\partial\phi},
  \qquad \phi = \frac{\pp}{J},
  \label{eq:m_mix}
\end{equation}
with the dimensionless interaction parameter
\begin{equation}
  \chi(\theta,\phi) = \chi_\infty + S_\chi\,\theta + \chi_1\,\phi,
  \label{eq:chi_nondim}
\end{equation}
where $\chi_\infty = \chi_0(T_\infty)$ and the thermo-swelling sensitivity is
\begin{equation}
  S_\chi = \Delta T_*\left.\frac{d\chi_0}{dT}\right|_{T_\infty}.
  \label{eq:Schi}
\end{equation}
For the linear concentration dependence in Eq.~\eqref{eq:chi_nondim},
$\partial\chi/\partial\phi=\chi_1$.

The elastic contribution follows from Eq.~\eqref{eq:mu_el}:
\begin{equation}
  m_\mathrm{el} = \Omega_e\,(J - J^{-1}),
  \qquad
  \Omega_e = \frac{G\,v_s}{R_g\,T_\infty},
\end{equation}
measuring the network stiffness relative to the thermal mixing energy.

\subsection{Dimensionless governing equations}

Dropping tildes for clarity, the dimensionless system reads:

\paragraph{Swelling equation.}
\begin{equation}
  \frac{\partial J}{\partial t}
  = \frac{\partial}{\partial x}\!\left(\Ms(J,\theta)\,\frac{\partial m}{\partial x}\right),
  \label{eq:J_nondim}
\end{equation}

\paragraph{Chemical potential.}
\begin{equation}
  m = m_\mathrm{mix}(J,\theta) + m_\mathrm{el}(J) - \ell^2\,\frac{\partial^2 J}{\partial x^2},
  \label{eq:m_nondim_box}
\end{equation}

\paragraph{Reactant equation.}
\begin{equation}
  \frac{\partial(J u)}{\partial t}
  + \frac{\partial}{\partial x}\!\left(u\,q - \delta\,\Ds_\mathrm{ref}(J,\theta)\,\frac{\partial u}{\partial x}\right)
  = -\Da\,J\,\Rhat(u,\theta,J),
  \label{eq:u_nondim}
\end{equation}
where $q = -\Ms(J,\theta)\,\partial_x m$ and
\begin{equation}
  \delta = \frac{D_0}{M_0\,\mu_*}
\end{equation}
is the ratio of solute diffusion speed to solvent permeation speed.

The Damk\"ohler number, measuring reaction rate relative to swelling rate, is
\begin{equation}
  \Da = \tau_s\,k_0\,c_b^{\,n-1}\,\exp\!\left(-\frac{E_a}{R_g T_\infty}\right).
  \label{eq:Da}
\end{equation}

\paragraph{Heat equation.}
\begin{multline}
  \Cs_\mathrm{ref}(J)\,\frac{\partial\theta}{\partial t}
  = \alpha\,\frac{\partial}{\partial x}\!\left(\Ks_\mathrm{ref}(J)\,\frac{\partial\theta}{\partial x}\right)
  + \Da\,J\,\Rhat(u,\theta,J)
  \\
  - \alpha\,\Pe_T\,\frac{\partial}{\partial x}\!\left[(\Theta_\infty + \theta)\,q\right],
  \label{eq:theta_nondim}
\end{multline}
with
\begin{equation}
  \alpha = \frac{K_*\,\tau_s}{C_*\,H_0^2} = \frac{\tau_s}{\tau_T},
  \quad
  \Pe_T = \frac{c_p^\ell\,M_0\,\mu_*}{K_*},
  \quad
  \Theta_\infty = \frac{T_\infty}{\Delta T_*},
  \label{eq:thermal_params}
\end{equation}
where $\tau_T = C_*H_0^2/K_*$ is the thermal diffusion time.
When $\Pe_T \ll 1$, the convective term may be dropped.

\paragraph{Constraint.}
\begin{equation}
  \phi = \frac{\pp}{J}.
\end{equation}

\subsection{Dimensionless boundary conditions}

At $x=0$ (symmetry):
\begin{equation}
  q = 0, \qquad
  u\,q - \delta\,\Ds_\mathrm{ref}\,u_x = 0, \qquad
  \theta_x = 0.
  \label{eq:bc0_nondim}
\end{equation}

At $x=1$ (free surface):
\begin{align}
  q &= \Bi_\mu\,(m - m_b), \label{eq:bc1_mu} \\
  u\,q - \delta\,\Ds_\mathrm{ref}\,u_x &= \Bi_c\,(u - 1), \label{eq:bc1_c} \\
  -\alpha\,\Ks_\mathrm{ref}(J)\,\theta_x &= \Bi_T\,\theta, \label{eq:bc1_T}
\end{align}
where the surface Biot numbers for permeation and mass transfer are
\begin{equation}
  \Bi_\mu = \frac{L_\mu\,H_0}{M_0},
  \qquad
  \Bi_c = \frac{k_c\,H_0}{M_0\,\mu_*},
  \label{eq:Bi_mu_c}
\end{equation}
and the thermal exchange number is
\begin{equation}
  \Bi_T = \frac{h_T\,\tau_s}{C_*\,H_0}
  = \alpha\,\frac{h_T\,H_0}{K_*},
  \label{eq:Bi_T}
\end{equation}
which measures the surface heat loss rate directly on the swelling
timescale $\tau_s$.  The standard thermal Biot number
$h_T H_0/K_*=\Bi_T/\alpha$ is recovered by factoring out $\alpha$.

\subsection{Controlling parameter groups}

The oscillatory dynamics are governed by the following dimensionless groups:
\begin{itemize}
  \item $\Da$\,---\,Damk\"ohler number: strength of reaction heat source relative to swelling.
  \item $\alpha = \tau_s/\tau_T$\,---\,thermal diffusion relative to swelling.
  \item $S_\chi$\,---\,thermo-swelling sensitivity: how strongly temperature shifts the interaction parameter.
  \item $\Gamma_A$\,---\,Arrhenius feedback strength.
  \item $\delta$\,---\,solute diffusion to permeation ratio.
  \item $\ell$\,---\,interface regularization length.
  \item $\Bi_c$\,---\,surface reactant Biot number.
  \item $\Bi_T$\,---\,thermal exchange number (surface heat loss on the swelling timescale).
\end{itemize}
Together with the working-point parameters $\pp$ and $\chi_\infty$, these form the minimal set controlling the onset of self-oscillation. In physical terms, oscillation requires that the fast positive feedback (reaction $\to$ heating $\to$ collapse) and the slow negative feedback (collapse $\to$ transport suppression $\to$ cooling $\to$ re-swelling) operate on well-separated time scales.

%% ================================================================
%%  Appendix C: Dispersion Relation Derivation
%% ================================================================

\end{document}
